\section{Results}
\label{sec:results}

All nine simulated pairs were complete. The method A and method B means were
\qty{5.12}{\milli\metre} and \qty{5.19}{\milli\metre}; the mean signed
difference was \qty{0.07}{\milli\metre} (\cref{tab:summary}).

\begin{table}[htbp]
  \centering
  \caption{Summary of the simulated paired values.}
  \label{tab:summary}
  \begin{tabular}{@{}lS[table-format=1.2]@{}}
    \toprule
    {Quantity} & {Value (\si{\milli\metre})} \\
    \midrule
    Method A mean & 5.12 \\
    Method B mean & 5.19 \\
    Mean difference & 0.07 \\
    \bottomrule
  \end{tabular}

  \smallskip\footnotesize\textit{Note.} Values are simulated and rounded for
  this typesetting demonstration.
\end{table}

\Cref{fig:paired-values} shows the paired values. The diagnostic panels in
\cref{fig:diagnostics} provide complementary views.

\begin{figure}[htbp]
  \centering
  \includegraphics[width=0.72\linewidth]{figures/method_comparison.pdf}
  \caption{Relationship between method A and method B for simulated paired
  values. The dashed line marks equality.}
  \label{fig:paired-values}
\end{figure}

\begin{figure}[htbp]
  \centering
  \begin{subfigure}[t]{0.48\textwidth}
    \centering
    \includegraphics[width=\linewidth]{figures/method_comparison.pdf}
    \caption{Values on the original scale.}
    \label{fig:diagnostic-values}
  \end{subfigure}\hfill
  \begin{subfigure}[t]{0.48\textwidth}
    \centering
    \includegraphics[width=\linewidth]{figures/difference_plot.pdf}
    \caption{Signed difference against pair mean.}
    \label{fig:diagnostic-differences}
  \end{subfigure}
  \caption{Two diagnostic views of the same simulated data.}
  \label{fig:diagnostics}
\end{figure}
